\section{The decomposition: depth reduction plus an index-coding crux}\label{sec:decomposition}

This section reduces the oblivious case of~(T) to two statements of different kinds. For an oblivious machine and a
fixed input length, the run is a circuit: the value computed at each step is a fixed function of the input bit under
the input head, of the previous state, and of the symbols under the work heads, and each of these symbols was written
at the previous visit of its head to its cell. Revealing the values at a set $\Gamma$ of steps cuts this circuit. If
every path of the circuit that avoids $\Gamma$ is short, each output bit depends only on the revealed values and on
few input bits, so the revealed values form an index code for the transpose demand with small side information.

\Cref{sec:dec-dag} sets up the circuit and proves the cone lemma. \Cref{sec:dec-D} states the two hypotheses, depth
reduction at linear cost (depth-TCP) and a lower bound for index codes on the resulting side information (the crux),
and shows that together they exclude oblivious linear-time transposers (Theorem~D). \Cref{sec:dec-rates} records what
the crux must deliver. \Cref{sec:dec-B} proves depth reduction for walks of linear leg area, in particular for
home-based walks, for every coupling of the heads (Theorem~B), and \cref{sec:dec-corB} turns its cuts into genies,
which reduces home-based oblivious machines to the crux (Corollary~B). None of the results of this section is formalized; each comes with a complete proof.

\subsection{The dependency DAG of an oblivious run}\label{sec:dec-dag}

Let $M$ be a machine as in \cref{sec:prelim-machines}, with state set $\{0,\dots,q\}$ and $k\ge1$ work tapes over the
alphabet $\{0,\dots,g\}$, and let $L\ge1$ be an input length at which $M$ is oblivious, with halting time~$T$. On
every input $w\in\{0,1\}^L$ the machine executes steps $0,\dots,T-1$. Write $x_0(t)$ for the position of the input
head and $x_j(t)$ for the position of work head~$j$ ($1\le j\le k$) at time~$t$; by obliviousness these positions,
and the set of steps that append an output bit, do not depend on~$w$. Step~$t$ reads the symbols under the heads at
time~$t$, writes the cells $x_1(t),\dots,x_k(t)$, moves the heads, and possibly appends an output bit. As in
\cref{sec:prelim-walks}, $\prev_j(t)=\max\{s<t:x_j(s)=x_j(t)\}$ when this set is nonempty.

The \emph{value} $v_t(w)$ of step~$t$ on input~$w$ consists of the state after step~$t$, the $k$ symbols written at
step~$t$, and the output bit appended at step~$t$ (or a symbol for ``none''). It takes at most $2^{\beta(M)}$ values,
where
\[
  \beta(M)=\clog{(q+1)}+k\,\clog{(g+1)}+2
\]
is the \emph{value width} of~$M$; the last two bits encode the output (none, $0$ or $1$).

\begin{definition}[Dependency DAG]\label{def:depgraph}
The \emph{dependency DAG} $H_M$ of $M$ at length~$L$ has vertex set $[0,T)=\{0,\dots,T-1\}$ and the following edges,
all pointing backwards in time: a \emph{state edge} $t\to t-1$ for every $1\le t<T$, and a \emph{cell edge}
$t\to\prev_j(t)$, labelled~$j$, for every work tape~$j$ and every $t$ for which $\prev_j(t)$ is defined. The
\emph{length} of an edge $t\to s$ is $t-s$.
\end{definition}

Every vertex has out-degree at most $k+1$. For each~$j$ the map $t\mapsto\prev_j(t)$ is injective: if
$\prev_j(t)=\prev_j(t')=s$ with $t<t'$, then $x_j(t')=x_j(s)=x_j(t)$ with $s<t<t'$, so $\prev_j(t')\ge t>s$. Hence
the symbol written on tape~$j$ at step~$s$ is read along at most one edge, the state after step~$s$ only along
$s+1\to s$, and the output bit along none. The DAG depends only on the head schedule of~$M$ at length~$L$, not on the
input or on the transition table.

\begin{lemma}[Local recursion]\label{lem:dec-local}
\notformalized
For every $t<T$ there is a function $\Delta_t$ such that for every input $w\in\{0,1\}^L$,
\[
  v_t(w)=\Delta_t\bigl(v_{t-1}(w),\,w_{x_0(t)},\,v_{\prev_1(t)}(w),\dots,v_{\prev_k(t)}(w)\bigr),
\]
where an undefined argument is replaced by a constant: by the initial state if $t=0$, by the end marker if
$x_0(t)\notin[0,L)$, and by the blank symbol if $\prev_j(t)$ is undefined.
\end{lemma}

\begin{informalproof}
Step~$t$ applies the transition function to the state at time~$t$, the input symbol under the input head and the
$k$ symbols under the work heads, and $v_t(w)$ is read off from its result. The state at time~$t$ is the initial
state if $t=0$ and the state component of $v_{t-1}(w)$ otherwise. The input symbol is $w_{x_0(t)}$, or the end marker
if $x_0(t)\notin[0,L)$. Every step writes the cell under each work head, so at time~$t$ the cell under work head~$j$
holds the symbol written on tape~$j$ at step $\prev_j(t)$, a component of $v_{\prev_j(t)}(w)$, and it is blank if
the head has not visited it before time~$t$; this is the cell-contents lemma of the first paper, Lean-verified
there~\cite[Lemma~2.3]{FirstPaper2026}. The positions $x_0(t)$ and the steps $\prev_j(t)$ do not depend on~$w$, so neither
does~$\Delta_t$.
\end{informalproof}

\begin{definition}[Genies, cones and depth]\label{def:genie}
A \emph{genie} is a set $\Gamma\subseteq[0,T)$ of revealed steps. For an integer $\beta\ge\beta(M)$, its
\emph{revealed word} $G_\Gamma(w)=(v_s(w))_{s\in\Gamma}$ is written with $\beta$ bits per value; this defines a map
$G_\Gamma\colon\{0,1\}^L\to\{0,1\}^{\beta\card\Gamma}$, and $\beta\card\Gamma$ is the \emph{cost} of~$\Gamma$. A path
$t_0\to t_1\to\dots\to t_m$ in $H_M$ has \emph{length}~$m$, and it is \emph{$\Gamma$-free} if
$t_1,\dots,t_m\notin\Gamma$ (its first vertex may lie in~$\Gamma$). The \emph{cone} $C_\Gamma(t)$ is the set of
endpoints of $\Gamma$-free paths starting at~$t$, so $t\in C_\Gamma(t)$. Its \emph{boundary} $\partial_\Gamma(t)$ is
the set of $s\in\Gamma$ such that some edge of $H_M$ leads from a vertex of $C_\Gamma(t)$ to~$s$, and
$J_\Gamma(t)=x_0(C_\Gamma(t))\cap[0,L)$ is the set of input positions read in the cone. The genie has \emph{depth at
most~$D$} if every $\Gamma$-free path has length at most~$D$.
\end{definition}

So revealing the value $v_s$ blocks every edge into~$s$.

\begin{lemma}[Cone lemma]\label{lem:cone}
\notformalized
Let $\Gamma$ be a genie and $t<T$.
\begin{enumerate}[label=(\alph*)]
\item If two inputs $w,w'\in\{0,1\}^L$ agree on $J_\Gamma(t)$ and satisfy $v_s(w)=v_s(w')$ for every
$s\in\partial_\Gamma(t)$, then $v_t(w)=v_t(w')$. Thus $v_t$ is a function of the values at the boundary and of the
input bits read in the cone.
\item If $\Gamma$ has depth at most~$D$, then
$\card{J_\Gamma(t)}\le\card{C_\Gamma(t)}\le\sum_{i=0}^{D}(k+1)^i\le(k+1)^{D+1}$.
\end{enumerate}
\end{lemma}

\begin{informalproof}
(a) By strong induction on~$t$. Let $w,w'$ be as stated; we show that the arguments of $\Delta_t$ in
\cref{lem:dec-local} agree on $w$ and~$w'$. Constant arguments agree trivially. If $x_0(t)\in[0,L)$, then
$x_0(t)\in J_\Gamma(t)$ because $t\in C_\Gamma(t)$, so $w_{x_0(t)}=w'_{x_0(t)}$. Now let $t\to t'$ be an edge, with
$t'=t-1$ or $t'=\prev_j(t)$. If $t'\in\Gamma$, then $t'\in\partial_\Gamma(t)$, since the edge starts at
$t\in C_\Gamma(t)$, and $v_{t'}(w)=v_{t'}(w')$ by hypothesis. If $t'\notin\Gamma$, then every $\Gamma$-free path
from~$t'$, preceded by the edge $t\to t'$, is a $\Gamma$-free path from~$t$. Hence $C_\Gamma(t')\subseteq C_\Gamma(t)$,
so $J_\Gamma(t')\subseteq J_\Gamma(t)$ and $\partial_\Gamma(t')\subseteq\partial_\Gamma(t)$, and the induction
hypothesis for $t'<t$ gives $v_{t'}(w)=v_{t'}(w')$. So all arguments agree, and $v_t(w)=v_t(w')$.

(b) Every vertex of $C_\Gamma(t)$ is the endpoint of a $\Gamma$-free path from~$t$ of length at most~$D$, and since
out-degrees are at most $k+1$, at most $(k+1)^i$ paths of length~$i$ start at~$t$. The last inequality holds because
$k\ge1$.
\end{informalproof}

\begin{corollary}[Output bits]\label{cor:dec-output}
\notformalized
Suppose that $M$ appends exactly $L$ output bits at length~$L$, the $\ell$-th one ($0\le\ell<L$) at step~$t_\ell$,
and let $\Gamma$ be a genie. Put $I_\Gamma(\ell)=J_\Gamma(t_\ell)$. If $G_\Gamma(w)=G_\Gamma(w')$ and $w,w'$ agree on
$I_\Gamma(\ell)$, then output bit~$\ell$ is the same on $w$ and on~$w'$. If $\Gamma$ has depth at most~$D$, then
$\card{I_\Gamma(\ell)}\le(k+1)^{D+1}$.
\end{corollary}

\begin{informalproof}
Output bit~$\ell$ is a component of $v_{t_\ell}$, and $\partial_\Gamma(t_\ell)\subseteq\Gamma$, so
$G_\Gamma(w)=G_\Gamma(w')$ gives $v_s(w)=v_s(w')$ for every $s\in\partial_\Gamma(t_\ell)$. Apply \cref{lem:cone}.
\end{informalproof}

\begin{remark}[Field-level genies]\label{rem:dec-field}
A genie may also reveal single \emph{fields}: the state after a step, or the symbol written on one tape. Each field
is read along at most one edge, so revealing it blocks exactly that edge, and the proof of \cref{lem:cone} applies
verbatim to paths that avoid the blocked edges. We do not need this refinement. Below, an edge $t\to s$ is blocked by
putting~$s$ into~$\Gamma$, at a cost of $\beta$ bits, which also bounds the width of a single field.
\end{remark}

\subsection{Two hypotheses and Theorem~D}\label{sec:dec-D}

The first hypothesis, depth-TCP (TCP stands for \emph{trajectory cut problem}), says that a genie of linear cost
reduces the depth of every dependency DAG of linear size to a constant.

\begin{definition}[depth-TCP]\label{def:depth-tcp}
Let $k\ge1$, $D\ge0$ and $\beta\ge2$ be integers and $C,\eps>0$ reals. The statement
$\mathrm{depth}\text{-}\TCP(k,C,\beta,\eps,D)$ says: for every machine~$M$ with $k$ work tapes and value width
$\beta(M)\le\beta$ there is~$d_0$ such that, for every $d\ge d_0$ at which $M$ is oblivious at length $L=d^2b(d)$
with halting time $T\le C(L+1)$, the dependency DAG $H_M$ at length~$L$ has a genie $\Gamma$ of cost
$\beta\card\Gamma\le\eps L$ and depth at most~$D$.
\end{definition}

This is a statement about head schedules: whether a genie of a given size and depth exists depends only on the DAG.
By the trade-off of \cref{prop:tradeoff} below, applied to a machine that copies its input to its output, depth-TCP
with $C\ge1$ can hold only if $\eps\ge\beta/(D+1)$.

For the second hypothesis recall the transpose map. Bit~$s$ of entry $(i,j)$ of a $d\times d$ matrix with $b$-bit
entries sits at input position $(di+j)b+s$, and in the encoding of the transpose at position
$\operatorname{tpos}((di+j)b+s)=(dj+i)b+s$~\cite[Section~2.2 and Lemma~2.5]{FirstPaper2026}. The map $\operatorname{tpos}$ is an involution of
$[0,L)$, so on input~$w$ the output bit of a transposer at position~$\ell$ is $w_{\operatorname{tpos}(\ell)}$.

\begin{definition}[Index codes and the crux]\label{def:crux}
Let $d\ge1$ and $L=d^2b(d)$. A \emph{side-information system} at size~$d$ is a family $I=(I(\ell))_{\ell\in[0,L)}$
of subsets of $[0,L)$: receiver~$\ell$ knows the input bits at the positions in $I(\ell)$ and demands
$w_{\operatorname{tpos}(\ell)}$. A map $G\colon\{0,1\}^L\to\{0,1\}^m$ is an \emph{index code} for~$I$ with $m$ bits
if for every~$\ell$ and all $w,w'\in\{0,1\}^L$ with $G(w)=G(w')$ that agree on $I(\ell)$ we have
$w_{\operatorname{tpos}(\ell)}=w'_{\operatorname{tpos}(\ell)}$.

For a class $\mathcal S$ of side-information systems and a real $\rho\ge0$, the statement $\Crux(\mathcal S,\rho)$
says: for every $\eta>0$ there is~$d_0$ such that every index code for every system of $\mathcal S$ at a size
$d\ge d_0$ has at least $(\rho-\eta)L$ bits. In other words, index codes for systems of~$\mathcal S$ have at least
$\rho L-o(L)$ bits, with the $o(L)$ uniform over~$\mathcal S$.

The class $\mathcal S_{k,C,\beta,D}$ consists of the systems $I_\Gamma=(I_\Gamma(\ell))_{\ell\in[0,L)}$ of
\cref{cor:dec-output}, where $d\ge1$; $M$ is a machine with $k$ work tapes and $\beta(M)\le\beta$ that is oblivious
at length $L=d^2b(d)$ with halting time $T\le C(L+1)$ and appends exactly $L$ output bits; and $\Gamma$ is a genie
of depth at most~$D$ in~$H_M$.
\end{definition}

The class $\mathcal S_{k,C,\beta,D}$ depends on each machine only through its head schedule and its output times; no
transition table enters. Its sets have at most $(k+1)^{D+1}$ elements, and it grows with~$D$. The crux is a statement
about \emph{all} index codes, linear or not, and it must be: the revealed word of a genie is an arbitrary function of
the input, determined by the transition table. A lower bound for linear index codes alone would cover only machines
whose revealed words are linear functions of the input, such as machines with an $\F_2$-linear data path.

\begin{remark}[Two readings of the crux]\label{rem:dec-crux}
(a) \emph{The crux restricts head schedules.} Call receiver~$\ell$ of a system~$I$ \emph{trivial} if
$\operatorname{tpos}(\ell)\in I(\ell)$, and let $\mathrm{triv}(I)$ be the number of trivial receivers. Broadcasting
the demanded bits of the non-trivial receivers is an index code with $L-\mathrm{triv}(I)$ bits, so
$\Crux(\mathcal S,\rho)$ implies $\mathrm{triv}(I)\le(1-\rho)L+o(L)$ for the systems $I\in\mathcal S$. For
$\mathcal S_{k,C,\beta,D}$ this says that at most $(1-\rho)L+o(L)$ output positions~$\ell$ have the property that the input
head reads position $\operatorname{tpos}(\ell)$ at the end of a $\Gamma$-free path that starts at the step $t_\ell$ appending
output bit~$\ell$.

(b) \emph{Storage codes.} For a system~$I$ let $N_I(p)=I(\operatorname{tpos}(p))$, the side information
of the receiver that demands $w_p$. If $G$ is an index code for~$I$ with $m$ bits, then every fibre $G^{-1}(z)$
is a binary storage code for $N_I$ (\cref{def:storage-code}): two words of the fibre that agree on $N_I(p)$
agree at~$p$. Some fibre has at least $2^{L-m}$ words. Hence, if every binary storage code for every $N_I$,
$I\in\mathcal S$, has at most $2^{(1-\rho)L+o(L)}$ words, uniformly, then $\Crux(\mathcal S,\rho)$ holds; in
particular, if every $N_I$ satisfies the rate bound~$c$ (\cref{def:polymatroid}), then $\Crux(\mathcal S,1-c)$ holds
by \cref{thm:entropic-rate}. Conversely, a linear storage code $U\subseteq\F_2^L$ for $N_I$ of dimension $\lambda L$
gives the linear index code $w\mapsto w+U\in\F_2^L/U$ with $(1-\lambda)L$ bits: if $w-w'\in U$ vanishes on
$N_I(p)$, it agrees there with $0\in U$, so it vanishes at~$p$. Restricted to linear codes, the crux at
rate~$\rho$ is therefore equivalent to the statement that binary linear storage codes for the $N_I$, $I\in\mathcal S$,
have dimension at most $(1-\rho)L+o(L)$, uniformly (the kernel of a linear index code is a linear storage code). For
general codes the two notions are related by the duality recalled in \cref{sec:prelim-codes}.
\end{remark}

\begin{theorem}[Theorem D: decomposition]\label{thm:D}
\notformalized
Let $k\ge1$, $D\ge0$ and $\beta\ge2$ be integers, $C>0$, and $0<\eps<\rho$. Suppose that the statements
$\mathrm{depth}\text{-}\TCP(k,C,\beta,\eps,D)$ and $\Crux(\mathcal S_{k,C,\beta,D},\rho)$ hold. Then for every
machine~$M$ with $k$ work tapes and value width at most~$\beta$ there is~$d_0$ such that for no $d\ge d_0$ is $M$
oblivious at length $L=d^2b(d)$ with halting time at most $C(L+1)$ and at the same time a transposer of all
$d\times d$ matrices with $b(d)$-bit entries. In particular, no oblivious machine with $k$ work tapes and value width
at most~$\beta$ transposes within $C(L+1)$ steps for all~$d$, or even for infinitely many~$d$.
\end{theorem}

\begin{informalproof}
Fix~$M$. Let $d_1$ be given by $\mathrm{depth}\text{-}\TCP(k,C,\beta,\eps,D)$ for~$M$, let $d_2$ be given by
$\Crux(\mathcal S_{k,C,\beta,D},\rho)$ for $\eta=(\rho-\eps)/2$, and put $d_0=\max(d_1,d_2)$. Suppose that for some
$d\ge d_0$ the machine is oblivious at length $L=d^2b(d)$ with halting time $T\le C(L+1)$ and transposes all
$d\times d$ matrices. Every word of length~$L$ encodes such a matrix (\cite[Lemma~2.5]{FirstPaper2026}, Lean-verified there), so on every input~$w$
the output has length~$L$ and its bit~$\ell$ is $w_{\operatorname{tpos}(\ell)}$; in particular, $M$ appends exactly
$L$ output bits. By depth-TCP there is a genie~$\Gamma$ with $\beta\card\Gamma\le\eps L$ and depth at most~$D$, so
$I_\Gamma\in\mathcal S_{k,C,\beta,D}$. By \cref{cor:dec-output}, if $G_\Gamma(w)=G_\Gamma(w')$ and $w,w'$ agree on
$I_\Gamma(\ell)$, then $w_{\operatorname{tpos}(\ell)}=w'_{\operatorname{tpos}(\ell)}$: the revealed word $G_\Gamma$ is
an index code for $I_\Gamma$ with $\beta\card\Gamma$ bits. The crux gives
$\beta\card\Gamma\ge(\rho-\eta)L=\tfrac12(\rho+\eps)L>\eps L$, a contradiction.
\end{informalproof}

\begin{corollary}[The two hypotheses refute oblivious (T)]\label{cor:dec-obl}
\notformalized
Suppose that for all integers $k\ge1$ and $\beta\ge2$ and every real $C>0$ there are reals $0<\eps<\rho$ and an
integer $D\ge0$ such that $\mathrm{depth}\text{-}\TCP(k,C,\beta,\eps,D)$ and $\Crux(\mathcal S_{k,C,\beta,D},\rho)$
hold. Then for every oblivious machine~$M$ and every constant~$c$ there is~$d_0$ such that for no $d\ge d_0$ does $M$
transpose all $d\times d$ matrices within $c(L+1)$ steps. In particular, no oblivious machine witnesses~(T).
\end{corollary}

\begin{informalproof}
If $c=0$ there is nothing to prove, since a machine that halts at time~$0$ writes no output. A machine without work
tapes first receives one idle work tape, which keeps it oblivious and changes neither its running time nor its output.
So let $M$ have $k\ge1$ work tapes. It is oblivious at every input length, and if it transposes all $d\times d$ matrices
within $c(L+1)$ steps, its halting time at length $L=d^2b(d)$ is at most $c(L+1)$. Apply \cref{thm:D} with $C=c$ and
$\beta=\beta(M)\ge2$.
\end{informalproof}

\paragraph{Why it matters.} Theorem~D isolates exactly two statements that together refute the oblivious case
of~(T): a structural statement about tape dependency DAGs, which mentions neither the data nor the transition table,
and an information-theoretic statement about index codes, which sees a machine only through its head schedule and
its output times. An impossibility proof for oblivious machines can therefore be split into a depth-reduction part
and a coding part; for
instance, a crux at rate~$\tfrac12$ for every depth~$D$ would combine with depth reduction at a single
$\eps<\tfrac12$. Theorem~B and Corollary~B below settle the depth-reduction part for home-based heads; in general
both parts are open (\cref{sec:depth,sec:codes}). On the construction side, Theorem~D says what an oblivious
transposer~$M$ with $k$ work tapes and running time $C(L+1)$ must do; let $\beta=\beta(M)$. At every size at which it
transposes and for all $\eps>0$ and $D\ge0$, either $H_M$ has no genie of cost at most $\eps L$ and depth at
most~$D$, or the system $I_\Gamma\in\mathcal S_{k,C,\beta,D}$ of such a genie has an index code with at most $\eps L$
bits, namely~$G_\Gamma$; by \cref{rem:dec-crux}(b), some fibre of $G_\Gamma$ is a storage code of rate at least
$1-\eps$ for $N_{I_\Gamma}$.

Theorem~D is an elementary composition of the cone lemma with definitions tailored to it, the multitape analogue of
the route from small-depth computation to index coding by depth reduction recalled in the introduction. Its content is
the formulation of the two hypotheses and the bookkeeping of genie cost against depth.

\begin{remark}[Relation to the first paper]\label{rem:dec-first}
The first paper proves, Lean-verified, that the plain form of its Tape Coding Conjecture, which forbids causal codes
for complete exchange on the block networks of tape diagrams of linear length, at the profile capacities, refutes the
oblivious case of~(T)~\cite[Theorem~9.9]{FirstPaper2026}. Theorem~D is not a refinement of that reduction, and we know no
implication, in either direction, between the plain Tape Coding Conjecture and the conjunction of depth-TCP and the crux:
depth-TCP concerns the step-level dependency DAG and says nothing about codes on block networks, and the crux
concerns index codes for cone side information rather than complete exchange.
\end{remark}

\subsection{What the crux must deliver}\label{sec:dec-rates}

Theorem~D needs the crux at a rate~$\rho$ above the relative cost~$\eps$ of the genie. The next proposition bounds that
cost from below in terms of the depth.

\begin{proposition}[Forced trade-off]\label{prop:tradeoff}
\notformalized
Let $M$ be oblivious at length~$L$ with halting time~$T$, and let $\Gamma$ be a genie of depth at most~$D$. Then
$\card\Gamma\ge\lfloor(T-1)/(D+1)\rfloor$. If moreover $T\ge L$, as for every machine that appends $L$ output bits
and in particular for every transposer, and $\beta\card\Gamma\le\eps L$, then
\[
  \eps\;\ge\;\frac{\beta}{D+1}-\frac{\beta}{L},\qquad\text{equivalently}\qquad D+1\;\ge\;\frac{\beta L}{\eps L+\beta}.
\]
\end{proposition}

\begin{informalproof}
Let $n=\lfloor(T-1)/(D+1)\rfloor$ and $a_i=i(D+1)$ for $0\le i<n$. Then $a_i+D+1\le n(D+1)\le T-1$, and the path
$a_i+D+1\to a_i+D\to\dots\to a_i$ of state edges has length $D+1$, so it is not $\Gamma$-free: one of
$a_i,\dots,a_i+D$ lies in~$\Gamma$. These $n$ windows are disjoint, so $\card\Gamma\ge n\ge(T-1-D)/(D+1)$. With
$T\ge L$ and $\beta\card\Gamma\le\eps L$ this gives $\eps L\ge\beta(L-D-1)/(D+1)=\beta L/(D+1)-\beta$, which is
either form of the claim.
\end{informalproof}

\paragraph{Why it matters.} In Theorem~D the crux is used at a rate $\rho>\eps$, and for a transposer
\cref{prop:tradeoff} gives $\eps\ge\beta/(D+1)-\beta/L$. So at depth~$D$ the crux must deliver a rate of at least
about $\beta/(D+1)$, and since $\beta\ge2$, genies of cost $\eps L$ have depth $D\ge2/\eps-1-o(1)$: the cheaper the
genie, the deeper the side information that the crux has to handle. The next proposition shows that the size of the
cones alone never suffices.

\begin{proposition}[The counting bound cannot replace the crux]\label{prop:dec-counting}
\notformalized
Let $I$ be a side-information system at size~$d$ with $\card{I(\ell)}\le B$ for all~$\ell$. Call a set~$W$ of
non-trivial receivers (\cref{rem:dec-crux}(a)) \emph{acyclic} if there are no distinct $\ell_1,\dots,\ell_r\in W$, $r\ge2$, with
$\operatorname{tpos}(\ell_{i+1})\in I(\ell_i)$ for all~$i$ (indices modulo~$r$), and let $\MAIS(I)$ be the largest
size of an acyclic set.
\begin{enumerate}[label=(\alph*)]
\item Every index code for~$I$ has at least $\MAIS(I)$ bits.
\item $\MAIS(I)\ge(L-\mathrm{triv}(I))/(B+1)$.
\item Suppose that a machine~$M$ with $k$ work tapes, oblivious at length $L$ with halting time $T\ge L$, has a genie
$\Gamma$ of depth at most~$D$ and cost $\beta\card\Gamma\le\eps L$, where $\beta\ge2$, $0<\eps\le1$ and
$L\ge\beta/\eps$. Then $(k+1)^{D+1}>1/\eps$. Consequently, for every system~$I$ at this size whose sets have at most
$B=(k+1)^{D+1}$ elements, in particular for every system of $\mathcal S_{k,C,\beta,D}$ (\cref{cor:dec-output}), the
lower bound $(L-\mathrm{triv}(I))/(B+1)$ given by (a) and~(b) is smaller than $\eps L$; it never contradicts the
existence of an index code with $\eps L$ bits, such as the revealed word $G_\Gamma$ when $M$ is a transposer.
\end{enumerate}
\end{proposition}

\begin{informalproof}
(a) This is the acyclic-subgraph bound of Bar-Yossef et al.~\cite[Theorem~3]{BBJK2011}; we include the short proof in
our notation. Let $W$ be acyclic, and order it as
$\ell_1,\dots,\ell_r$ so that $\operatorname{tpos}(\ell_j)\in I(\ell_i)$ with $j\ne i$ implies $j<i$; this is
possible because the digraph on~$W$ with an arc from $\ell$ to $\ell'$ when $\operatorname{tpos}(\ell')\in I(\ell)$
has no cycle. Let $Z$ be the set of inputs that vanish outside $\operatorname{tpos}(W)$; it has $2^{\card W}$
elements because $\operatorname{tpos}$ is injective. We show that an index code~$G$ is injective on~$Z$. Let
$w,w'\in Z$ with $G(w)=G(w')$, and suppose $w$ and $w'$ agree at $\operatorname{tpos}(\ell_j)$ for all $j<i$. A
position in $I(\ell_i)$ lies outside $\operatorname{tpos}(W)$, where both inputs vanish, or equals
$\operatorname{tpos}(\ell_j)$ for some $j<i$, because $\operatorname{tpos}(\ell_i)\notin I(\ell_i)$ (receiver
$\ell_i$ is non-trivial). So $w$ and $w'$ agree on $I(\ell_i)$, hence at $\operatorname{tpos}(\ell_i)$. By induction
$w=w'$, so $2^m\ge2^{\card W}$.

(b) Order the non-trivial receivers uniformly at random, and let $W$ consist of those receivers~$\ell$ that precede
every non-trivial $\ell'\ne\ell$ with $\operatorname{tpos}(\ell')\in I(\ell)$. Then $W$ is acyclic, since along a
cycle in~$W$ each receiver would precede the next. As $\operatorname{tpos}$ is injective, a receiver~$\ell$ has at
most $\card{I(\ell)}\le B$ such receivers~$\ell'$, so it lies in~$W$ with probability at least $1/(B+1)$, and
$\mathbb E\card W\ge(L-\mathrm{triv}(I))/(B+1)$.

(c) By \cref{prop:tradeoff}, $D+1\ge\beta L/(\eps L+\beta)\ge\beta L/(2\eps L)=\beta/(2\eps)\ge1/\eps$, using
$\beta\le\eps L$ and $\beta\ge2$. Hence $(k+1)^{D+1}\ge2^{D+1}\ge2^{1/\eps}>1/\eps$, and
$(L-\mathrm{triv}(I))/(B+1)\le L/(B+1)<\eps L$.
\end{informalproof}

\paragraph{Why it matters.} The combined bound $(L-\mathrm{triv}(I))/(B+1)$ of (a) and~(b) uses only the sizes of the
sets of the side information and the number of trivial receivers. Part~(c) shows that it is always below $\eps L$: the
depth that a genie of cost $\eps L$ must reach is so large that the size bound $(k+1)^{D+1}$ on the cones is far bigger
than $1/\eps$. No argument that uses only the sizes can do better, because splitting the receivers into groups of
$B+1$, each receiver knowing the bits demanded by the others in its group, gives a system with sets of size~$B$ and an
index code of $\lceil L/(B+1)\rceil$ bits, one XOR per group. A proof of the crux must use the structure of the side
information produced by tape schedules.

\begin{remark}[The rate needed]\label{rem:dec-rates}
For each depth~$D$, let $\rho^*(D)$ be the supremum of all rates~$\rho$ such that the statement
$\Crux(\mathcal S_{k,C,\beta,D},\rho)$ holds. It is non-increasing in~$D$, because the classes grow with~$D$. Suppose that depth-TCP holds with parameters
$(k,C,\beta,\eps,D(\eps))$ for all small $\eps>0$. By Theorem~D, oblivious machines with these parameters are then
excluded as soon as $\rho^*(D(\eps))>\eps$ for a single~$\eps$. If $D(\eps)\le c\,\eps^{-a}$ for constants $c$
and~$a$ (by \cref{prop:tradeoff}, necessarily $a\ge1$ if depth-TCP is to apply to transposers), a crux with
$\rho^*(D)\ge D^{-1/a+\delta}$ for some $0<\delta<1/a$ suffices, because then
\[
  \rho^*(D(\eps))\;\ge\;\bigl(c\,\eps^{-a}\bigr)^{-1/a+\delta}\;=\;c^{-1/a+\delta}\,\eps^{1-a\delta}\;>\;\eps
\]
for all small~$\eps$. For heads of leg area at most $AT$ (home-based heads have $A=1$), Corollary~B below gives depth
$D(\eps)=2^{O(k^3(1+A)\beta C/\eps)}$; there a crux at a single depth~$D$ with a rate exceeding about
$456k^3(1+A)\beta C/\log_2D$ suffices. Both requirements are far
above the rate, at most $(k+1)^{-D-1}$, that the counting bound of \cref{prop:dec-counting} gives.
\end{remark}

\subsection{Depth reduction for home-based walks: Theorem~B}\label{sec:dec-B}

Theorem~B is a statement about walks; machines reappear in \cref{sec:dec-corB}. Throughout this subsection
$x_1,\dots,x_k$ are walks of length~$T$ (\cref{sec:prelim-walks}), indexed by \emph{heads} $h\in\{1,\dots,k\}$. They
may be coupled arbitrarily: no relation between them is assumed.

\paragraph{The walk DAG and cuts.} The \emph{walk DAG} of $x_1,\dots,x_k$ has vertex set $[0,T)$ and two kinds of
edges: \emph{state edges} $t\to t-1$ for $1\le t<T$, and \emph{long hops} $u\to\prev_h(u)$, one for every head~$h$
and every time~$u$ with $u-\prev_h(u)\ge2$ (the hop $(\prev_h(u),u)$ of \cref{sec:prelim-walks}, read backwards). A
pause $u\to u-1$ of a head joins the same two vertices as the state edge $u\to u-1$ and is identified with it; in
general an edge is determined by its two endpoints, so a pair that is a long hop of several heads is a single edge. A
\emph{cut} is a set~$F$ of edges, and it has \emph{depth at most~$D$} if every path that uses no edge of~$F$ contains
at most~$D$ long hops. For the work-head walks of an oblivious machine~$M$ at length~$L$, the walk DAG has the
vertices of $H_M$ (\cref{def:depgraph}), and every edge of $H_M$ joins the endpoints of an edge of the walk DAG: a
cell edge of length~$1$ is a pause, and a cell edge of length at least~$2$ is a long hop.

\paragraph{Bands, legs and leg area.} For $z\in\Z$ let $\lambda(z)=\lfloor\log_2\abs z\rfloor$ if $\abs z\ge2$ and
$\lambda(z)=0$ otherwise. The \emph{band} of head~$h$ at time~$t$ is $\lambda_h(t)=\lambda(x_h(t))$, a function of
the cell. Bands, scales, floors and the exponents~$\theta$ of radii $2^\theta$ are integers throughout. For integers
$\theta\ge1$, $\lambda(z)\ge\theta$ if and only if $\abs z\ge2^\theta$; and since $\abs{x_h}$ changes by at most~$1$ per
step, so does~$\lambda_h$. The \emph{scale} of a leg (\cref{sec:prelim-walks}) is the largest band over it. For
$\theta\ge0$ let $N_h(\theta,t)$ be the number of legs of head~$h$ of scale at least~$\theta$ that begin at or before
time~$t$, and let $N_h(\theta)=N_h(\theta,T-1)$. The \emph{leg area} of head~$h$ is
\[
  S_h=\sum_{\text{legs $\ell$ of $h$}}2^{\mathrm{scale}(\ell)},
\]
and the \emph{area hypothesis} with constant $A>0$ is
\[
  \text{(A)}\qquad S_h\le A\,T\quad\text{for every head }h.
\]
Every leg contributes at least~$1$, so head~$h$ has at most $S_h$ legs. The statements below use only that the legs
of each walk partition $[0,T)$ into intervals on which the walk is weakly monotone.

\begin{lemma}[Leg area of home-based walks]\label{lem:dec-area}
\notformalized
If every leg of a walk~$x$ of length~$T$ begins or ends at a cell~$z$ with $\abs z\le1$, then every leg~$\ell$
satisfies $2^{\mathrm{scale}(\ell)}\le\card\ell$, and the leg area of~$x$ is at most~$T$. In particular, home-based
walks satisfy~(A) with $A=1$.
\end{lemma}

\begin{informalproof}
If the scale~$s$ of $\ell$ is~$0$, then $2^s=1\le\card\ell$. Otherwise $\ell$ contains a time~$t$ with
$\abs{x(t)}\ge2^s$ and has an endpoint~$e$ with $\abs{x(e)}\le1$. Since $x$ moves by at most~$1$ per step,
$\abs{t-e}\ge2^s-1$, and the interval~$\ell$ contains the $\abs{t-e}+1\ge2^s$ times between $e$ and~$t$. Summing
over the legs, which partition $[0,T)$, bounds the leg area by~$T$.
\end{informalproof}

The walk DAG does not change when a single walk is translated, so everything below applies to heads with individual
fixed home cells, the bands of each head being measured from its home.

\begin{lemma}[Crossings and spanning hops]\label{lem:dec-facts}
\notformalized
Let $x$ be a walk of length~$T$ with leg area~$S$, let $\theta\ge1$, and let $N(\theta)$ be the number of legs of~$x$
of scale at least~$\theta$. A \emph{crossing} of radius~$2^\theta$ at time~$t$ ($1\le t<T$) is \emph{outward} if
$\abs{x(t-1)}<2^\theta\le\abs{x(t)}$ and \emph{inward} if $\abs{x(t-1)}\ge2^\theta>\abs{x(t)}$.
\begin{enumerate}[label=\textup{(F\arabic*)}]
\item The walk has at most $2N(\theta)$ outward crossings, and at most $5N(\theta)$ crossings in total, of
radius~$2^\theta$. Moreover $\sum_{\theta\ge1}2^\theta N(\theta)\le2S$.
\item For every time~$\tau$ and every integer $R\ge1$, at most $2R-1$ hops $u\to\prev(u)$ of~$x$ \emph{span}~$\tau$,
that is, satisfy $\prev(u)<\tau\le u$, and have $\abs{x(u)}<R$.
\end{enumerate}
\end{lemma}

\begin{informalproof}
(F1) Assign each outward crossing~$t$ to the leg~$\ell$ that contains~$t$; the scale of~$\ell$ is at least~$\theta$
because $\abs{x(t)}\ge2^\theta$. Either $t$ is the first time of~$\ell$, which happens for at most one crossing per
leg, or $t-1\in\ell$ as well. In the second case suppose that $x$ is non-decreasing on~$\ell$; the non-increasing case
follows by the symmetry $z\mapsto-z$. If $x(t)\le0$, then $x(t-1)\le x(t)\le0$ and $\abs{x(t-1)}\ge\abs{x(t)}$,
which is impossible; so $x(t)\ge2^\theta$, and $x\ge2^\theta$ at all later times of~$\ell$, so no later outward
crossing $t'$ with $t'-1\in\ell$ is assigned to~$\ell$. Thus each leg of scale at least~$\theta$ receives at most two
outward crossings and every other leg none, which gives the first bound. Crossings of radius~$2^\theta$ alternate
between outward and inward, as each changes whether $\abs x\ge2^\theta$. So the inward crossings number at most the
outward ones plus one, and the extra one occurs only if $\abs{x(0)}\ge2^\theta$, in which case the first leg has
scale at least~$\theta$ and $N(\theta)\ge1$. Altogether there are at most $2N(\theta)+3N(\theta)$ crossings. Finally,
$\sum_{\theta\ge1}2^\theta N(\theta)=\sum_\ell\sum_{1\le\theta\le\mathrm{scale}(\ell)}2^\theta
\le\sum_\ell2^{\mathrm{scale}(\ell)+1}=2S$.

(F2) For such a hop, $x$ visits the cell $x(u)$ at time $\prev(u)<\tau$ and at no time strictly between $\prev(u)$
and~$u$, so $u$ is the first visit to $x(u)$ at or after~$\tau$. Hence each of the $2R-1$ cells~$z$ with $\abs z<R$
is the cell of at most one such hop.
\end{informalproof}

\begin{lemma}[Legs and long hops]\label{lem:legs}
\notformalized
Let $u\to p$ be a long hop of head~$h$, so $p=\prev_h(u)\le u-2$. The leg of~$h$ that contains~$u$ begins after~$p$
and has scale at least $\lambda_h(u)$. Hence $N_h(\theta,u)\ge N_h(\theta,p)+1$ for every $\theta\le\lambda_h(u)$.
\end{lemma}

\begin{informalproof}
Pick~$t$ with $p<t<u$. By the definition of $\prev_h$, head~$h$ is not at the cell $x_h(u)=x_h(p)$ at time~$t$. If
$p$ and~$u$ lay in one leg, $x_h$ would be weakly monotone on an interval containing $p<t<u$, so $x_h(t)$ would lie
between the equal values $x_h(p)$ and $x_h(u)$, a contradiction. So the leg containing~$u$ is a later interval than
the leg containing~$p$, and it begins after~$p$; its scale is at least $\lambda_h(u)$ because it contains~$u$. For
$\theta\le\lambda_h(u)$ this leg is counted in $N_h(\theta,u)$ but not in $N_h(\theta,p)$, and every leg counted in
$N_h(\theta,p)$ is counted in $N_h(\theta,u)$.
\end{informalproof}

\Cref{lem:legs} fails for pauses, whose two endpoints may lie in one leg; this is why pauses are treated as state edges.

\begin{definition}[The band label and its cut]\label{def:dec-label}
Fix integers $a\ge1$ (the \emph{group width}), $r\in\{0,\dots,a-1\}$ (the \emph{offset}), $K\ge2$ (the
\emph{modulus}) and $\omega_1,\dots,\omega_k\in\{0,\dots,K-1\}$.
\begin{itemize}[leftmargin=2em]
\item The \emph{group} of head~$h$ at time~$t$ is $g_h(t)=\lfloor(\lambda_h(t)+r)/a\rfloor$, and the \emph{floor} of
a group $G\ge0$ is $\theta(G)=\max(aG-r,0)$, the smallest band in group~$G$. The floors of the groups $G\ge1$ are the
integers $\theta\ge1$ with $\theta\equiv-r\pmod a$.
\item The \emph{occupied groups} at time~$t$ are the distinct values among $g_1(t),\dots,g_k(t)$, listed from the top
as $G_1(t)>G_2(t)>\dots>G_{m(t)}(t)$, with floors $\theta_i(t)=\theta(G_i(t))$; here $m(t)\le k$.
\item The \emph{shared counts} are $n_i(t)=\sum_{h=1}^kN_h(\theta_i(t),t)$ for $i\le m(t)$: the number of
legs, of all heads, of scale at least $\theta_i(t)$ begun by time~$t$.
\item The \emph{label} is
\[
  \Psi(t)=\sum_{i=1}^{m(t)}\bigl((n_i(t)+\omega_i)\bmod K\bigr)\,K^{k-i}\;\in[0,K^k).
\]
\item The \emph{cut} $F=F(a,r,K,\omega)$ consists of the long hops $u\to p$ with $\Psi(p)\ge\Psi(u)$ and the state
edges $t\to t-1$ with $\Psi(t-1)>\Psi(t)$.
\end{itemize}
\end{definition}

So $\Psi(t)$ is the integer whose base-$K$ digits, most significant first, are $\psi_i(t)=(n_i(t)+\omega_i)\bmod K$
for $i\le m(t)$ and $0$ for $m(t)<i\le k$, and labels compare as these digit strings compare lexicographically.
Whether an edge is cut depends only on its endpoints. For a fixed~$\theta$, the shared count $\sum_hN_h(\theta,t)$ is
non-decreasing in~$t$; this monotonicity and \cref{lem:legs} are what the label uses.

\begin{theorem}[Theorem B: depth reduction for home-based walks]\label{thm:B}
\notformalized
Let $x_1,\dots,x_k$ be walks of length~$T$, coupled arbitrarily, that satisfy the area hypothesis~(A) with
constant~$A$. Let $a\ge1$ and $K\ge2$ be integers, and let $r\in\{0,\dots,a-1\}$ and $\omega\in\{0,\dots,K-1\}^k$ be
uniformly random and independent. Then the cut $F=F(a,r,K,\omega)$ has depth at most $K^k-1$, and
\begin{equation}\label{eq:dec-B}
  \mathbb E_{r,\omega}\card F\;\le\;\frac{(20k^2A+k)\,T}{a}+\Bigl(9+\frac{2^{a+2}}a\Bigr)\frac{k^2AT}{K}
  \;\le\;20k^2(1+A)\,T\Bigl(\frac1a+\frac{2^a}{aK}\Bigr)+\frac{20k^2AT}{K}.
\end{equation}
Consequently:
\begin{enumerate}[label=(\alph*)]
\item for every $D\ge4^k$, the choice $K=\lfloor D^{1/k}\rfloor$ and $a=\lfloor\log_2D/(2k)\rfloor$, with suitable
offsets, gives a cut of depth at most $D-1$ with $\card F\le114k^3(1+A)T/\log_2D$;
\item for every $0<\eps\le1$, some cut of size at most $\eps T$ has depth less than
$2^{\lceil114k^3(1+A)/\eps\rceil}$;
\item in particular, for every coupling of $k$ home-based walks, $\eps T$ edges suffice to reach depth less than
$2^{\lceil228k^3/\eps\rceil}=2^{O(k^3/\eps)}$.
\end{enumerate}
\end{theorem}

\begin{informalproof}
\emph{Depth.} Along a path that avoids~$F$, $\Psi$ does not increase along state edges and strictly decreases along
long hops. Since $\Psi$ takes values in $[0,K^k)$, the path contains at most $K^k-1$ long hops.

\emph{Notation.} Fix the walks, $a$ and~$K$; probabilities and expectations are over $r$ and~$\omega$. We bound the
number of cut long hops counted with their head, as pairs $(h,u)$, which is at least the number of cut long-hop
edges, and the number of cut state edges. Consider a long hop $u\to p$ of head~$h$, with cell $z=x_h(u)=x_h(p)$.
Since bands are functions of the cell, $h$ has the same band $\lambda=\lambda(z)$ and the same group
$G=g_h(u)=g_h(p)$ at $p$ and at~$u$. Let $\theta^+=\theta(G+1)=a(G+1)-r\ge1$ be the floor of the next group: a head
is in a group above~$G$ at time~$t$ exactly when its cell has absolute value at least $2^{\theta^+}$, and
$\abs z<2^{\theta^+}$. Let $U(t)$ be the set of occupied groups above~$G$ at time~$t$. Whether $U(p)=U(u)$ holds,
and all counts $n_i$, depend on~$r$ only.

\emph{Step 1: $U(p)=U(u)$.} Then $G$ has the same position $j=\card{U(p)}+1$ in the lists at $p$ and at~$u$, and
$G_i(p)=G_i(u)$, with a common floor~$\theta_i$, for every $i\le j$. Let $\Delta_i=n_i(u)-n_i(p)$ for $i\le j$. For
$i<j$, $\Delta_i\ge0$ by monotonicity at a fixed floor. For $i=j$ we have $\theta_j=\theta(G)\le\lambda=\lambda_h(u)$,
so \cref{lem:legs} gives $N_h(\theta_j,u)\ge N_h(\theta_j,p)+1$, and with monotonicity for the other heads, $\Delta_j\ge1$.
Say that digit~$i$ \emph{wraps} if $\lfloor(n_i(p)+\omega_i)/K\rfloor\ne\lfloor(n_i(u)+\omega_i)/K\rfloor$. If no
digit $i\le j$ wraps, then $\psi_i(p)=\psi_i(u)-\Delta_i$ for every $i\le j$, so the first position at which the
digit strings of $p$ and~$u$ differ is at most~$j$ and carries the smaller digit at~$p$. Thus $\Psi(p)<\Psi(u)$, and
the hop is not cut; the digits below position~$j$ play no role. Given~$r$, since $\omega_i$ is uniform on
$\{0,\dots,K-1\}$, digit~$i$ wraps with probability at most $\Delta_i/K$. So, given~$r$, a hop of Step~1 is cut with
probability at most $\sum_{i\le j}\Delta_i/K$.

\emph{Step 2: $U(p)\ne U(u)$.} We count every such hop as cut and charge it to a crossing. Let $G'>G$ be a group that
lies in exactly one of $U(p)$ and $U(u)$. Since $h$ is in group~$G$ at both times, at one of the two times $G'$ is
occupied by a head $h'\ne h$ that is not in~$G'$ at the other time. Groups change by at most one per step, and a
head moves from a group $G''$ to $G''+1$ exactly when its band reaches $\theta(G''+1)$, that is, at an outward
crossing of radius $2^{\theta(G''+1)}$; it moves back at an inward crossing of the same radius. If $h'$ is in~$G'$
at time~$p$, consider its first move out of~$G'$ after~$p$; if $h'$ is in~$G'$ at time~$u$, consider its last move
into~$G'$ at or before~$u$. In either case the move happens at a step $\tau-1\to\tau$ with $p<\tau\le u$, and it is a
crossing of radius $2^{\theta(G'+1)}$ or $2^{\theta(G')}$. So $h'$ has a crossing of radius $2^{\theta^*}$ at a time
$\tau\in(p,u]$, where $\theta^*$ is a floor with $\theta^*\ge\theta(G')\ge\theta^+$.
Charge the hop to one such crossing $(h',\theta^*,\tau)$, chosen by a fixed rule. Every hop $(h,u)$ charged to
$(h',\theta^*,\tau)$ satisfies $\prev_h(u)<\tau\le u$ and $\abs{x_h(u)}=\abs z<2^{\theta^+}\le2^{\theta^*}$, so by
(F2) with $R=2^{\theta^*}$, applied to each of the $k$ heads, at most $k\,2^{\theta^*+1}$ hops are charged to each
crossing. By (F1), head~$h'$ has at most $5N_{h'}(\theta)$ crossings of radius~$2^\theta$. Hence, for fixed~$r$,
\[
  \#\{\text{hops of Step 2}\}\;\le\;\sum_{h'=1}^k\ \sum_{\theta\ge1\text{ a floor}}5N_{h'}(\theta)\cdot k\,2^{\theta+1}.
\]
Each $\theta\ge1$ is a floor for exactly one value of~$r$, hence with probability $1/a$, and by (F1) and~(A),
\[
  \mathbb E\,\#\{\text{hops of Step 2}\}\;\le\;\frac{10k}a\sum_{h'}\sum_{\theta\ge1}2^\theta N_{h'}(\theta)
  \;\le\;\frac{10k}a\sum_{h'}2S_{h'}\;\le\;\frac{20k^2AT}a.
\]

\emph{Step 3: wraps.} By Step~1, the expected number of cut hops of Step~1 is at most
$\frac1K\mathbb E\sum\sum_{i\le j}\Delta_i$, where the outer sum runs over the hops of Step~1. For such a hop and
$i\le j$, $\Delta_i$ is the number of legs, of all heads, that have scale at least~$\theta_i$ and begin in $(p,u]$.
We exchange the order of summation: for each leg~$\ell$, of any head, beginning at time~$\tau_\ell$ with
scale~$s_\ell$, we count the hops that span~$\tau_\ell$ in the sense of (F2), that is, satisfy $p<\tau_\ell\le u$. By the argument
of (F2), for each head and each cell at most one hop of that head spans~$\tau_\ell$ and ends at that cell.

\emph{Own digit ($i=j$).} A hop spanning~$\tau_\ell$ contributes~$\ell$ to its $\Delta_j$ only if $\theta(G)\le
s_\ell$, where $G$ is its group. Since $\theta(G)$ is the largest floor at most~$\lambda$ (or~$0$ if there is none),
this holds if $\lambda\le s_\ell$, and for $\lambda>s_\ell$ it holds exactly when no floor lies in $(s_\ell,\lambda]$.
The floors at least~$1$ form an arithmetic progression of difference~$a$ with a uniformly random offset, so the latter
has probability $\max(0,1-(\lambda-s_\ell)/a)$. There are $2^{s+2}-1$ cells of band at most~$s$, and $2^{\lambda+1}$
cells of band $\lambda\ge1$. So the expected number of hops counted for~$\ell$ is at most
\[
  k\Bigl(2^{s_\ell+2}+\sum_{i=1}^{a-1}2^{s_\ell+i+1}\Bigl(1-\frac ia\Bigr)\Bigr)
  \;\le\;k\Bigl(2^{s_\ell+2}+\frac{2^{s_\ell+a+2}}a\Bigr),
\]
because $\sum_{i=1}^{a-1}2^i(a-i)=2^a\sum_{j=1}^{a-1}j2^{-j}<2^{a+1}$. Summing over the legs of all heads, with
$\sum_\ell2^{s_\ell}=\sum_{h'}S_{h'}\le kAT$, gives $\mathbb E\sum\Delta_j\le4k^2AT(1+2^a/a)$.

\emph{Higher digits ($i<j$).} Here $G_i>G$, so $\theta_i\ge\theta^+$ is a floor at least~$1$ that exceeds the band
of the hop, and distinct~$i$ give distinct floors. So, for each leg~$\ell$, the pairs consisting of a hop
spanning~$\tau_\ell$ and an index $i<j$ with $\theta_i\le s_\ell$ number at most the sum, over the floors
$1\le\theta\le s_\ell$, of the number of hops spanning~$\tau_\ell$ whose cell has absolute value less than~$2^\theta$.
By (F2) this number is at most $k(2^{\theta+1}-1)$, and since consecutive floors differ by $a\ge1$, the sum is at most
$k\sum_{\theta=1}^{s_\ell}2^{\theta+1}<k\,2^{s_\ell+2}$. Summing over the legs gives
$\sum\sum_{i<j}\Delta_i\le4k^2AT$ for every~$r$.

Together, the expected number of cut hops of Step~1 is at most
$(4k^2AT(1+2^a/a)+4k^2AT)/K=(8+2^{a+2}/a)\,k^2AT/K$.

\emph{Step 4: state edges.} Consider the state edge $t\to t-1$. If no head changes group at step $t-1\to t$, the
occupied groups and their floors are the same at $t-1$ and~$t$, and $n_i(t-1)\le n_i(t)$ for every~$i$; as in
Step~1, the edge is cut only if some digit wraps, which, given~$r$, has probability at most
$\sum_i(n_i(t)-n_i(t-1))/K$. Here $n_i(t)-n_i(t-1)$ is the number of heads that begin, at time~$t$, a leg of scale at
least~$\theta_i$. A head begins at most one leg at a time and there are at most $k$ digits, so the sum over $i$
and~$t$ is at most $k$ times the number of legs, which is at most $\sum_{h'}S_{h'}\le kAT$; so these cuts number at
most $k^2AT/K$ in expectation. Otherwise some head changes group at step $t-1\to t$, and we count the edge as cut. A
head changes group at a step only if its band moves between two values $\lambda$ and $\lambda+1$ and $\lambda+1$ is a
floor, which has probability $1/a$; so the expected number of such steps is at most $kT/a$.

\emph{Total.} Adding Steps 2, 3 and~4,
\[
  \mathbb E\card F\;\le\;\frac{20k^2AT}a+\Bigl(8+\frac{2^{a+2}}a\Bigr)\frac{k^2AT}K+\frac{k^2AT}K+\frac{kT}a,
\]
which is the first bound of \eqref{eq:dec-B}; the second follows from $k\le20k^2$, $9\le20$ and
$2^{a+2}A\le20(1+A)2^a$. Some offsets $r,\omega$ achieve at most the expectation.

\emph{Parameters.} For~(a) put $y=\log_2D/(2k)\ge1$. Then $a=\lfloor y\rfloor\ge y/2$, and $K=\lfloor2^{2y}\rfloor
\ge2^{2y-1}$ because $2^{2y}\ge4$; moreover $K^k\le D$. Using $2^{-y}\le1/y$ for $y>0$,
\[
  \frac1a\le\frac2y=\frac{4k}{\log_2D},\qquad\frac{2^a}K\le2^{1-y}\le\frac2y=\frac{4k}{\log_2D},\qquad
  \frac1K\le2^{1-2y}\le2^{-y}\le\frac1y=\frac{2k}{\log_2D}.
\]
Inserting these into the first bound of \eqref{eq:dec-B}, with $2^{a+2}/(aK)\le4\cdot2^a/K$, gives
\[
  \mathbb E\card F\;\le\;\bigl(80k^3A+4k^2+18k^3A+16k^3A\bigr)\frac T{\log_2D}\;\le\;\frac{114k^3(1+A)\,T}{\log_2D},
\]
and the depth is at most $K^k-1\le D-1$. For~(b) take $\log_2D=\lceil114k^3(1+A)/\eps\rceil$, which is at least~$2k$.
Part~(c) is~(b) with $A=1$, by \cref{lem:dec-area}.
\end{informalproof}

\paragraph{Why it matters.} Theorem~B proves the walk form of the first hypothesis of Theorem~D for all heads of
linear leg area, in particular for home-based heads, and for every coupling: it makes no assumption on hop lengths, on
independence, or on how the heads move relative to each other. So, for an oblivious machine running in time at most
$C(L+1)$, a genie of cost $\eps L$ and depth at most $D(\eps)$, depending on $k$, $A$, $\beta$ and~$C$, exists for every
$\eps\in(0,1]$ (see the paragraph after the proof of \cref{cor:B}) unless the leg area of some work head is not $O(T)$
for any choice of home cell. That requires many legs whose duration is much shorter than
$2^{\mathrm{scale}}$; by \cref{lem:dec-area} such legs neither begin nor end near the home cell, so the head turns
around far from every fixed home. \Cref{sec:depth} studies walks of this kind.

\subsection{Home-based oblivious machines reduce to the crux}\label{sec:dec-corB}

Theorem~B cuts edges, counts only long hops, and pays one unit per cut edge. Two conversions turn its cut into a genie
in the sense of \cref{def:genie}: an edge $t\to s$ is blocked by revealing~$s$, at a cost of $\beta$ bits, and the
runs of short edges are cut by revealing periodic steps.

\begin{lemma}[Long-edge conversion]\label{lem:dec-longedge}
\notformalized
Let $M$ be oblivious at length~$L$ with halting time~$T$, let $\mu\ge1$ and $P>\mu$ be integers, and let $\Gamma_0$
be a genie such that every $\Gamma_0$-free path has at most $D'$ edges of length greater than~$\mu$. Then the genie
\[
  \Gamma=\Gamma_0\cup\{t\in[0,T):t\bmod P\ge P-\mu\}
\]
has depth at most $D'+(D'+1)(P-\mu)\le(D'+1)P-1$, and $\card\Gamma\le\card{\Gamma_0}+\mu\lceil T/P\rceil$; for
$\mu=1$, $\card\Gamma\le\card{\Gamma_0}+\lfloor T/P\rfloor$.
\end{lemma}

\begin{informalproof}
Let $t_0\to t_1\to\dots\to t_m$ be $\Gamma$-free. It is $\Gamma_0$-free, so at most $D'$ of its edges have length
greater than~$\mu$, and removing them leaves at most $D'+1$ maximal runs of consecutive edges of length at most~$\mu$.
Let $t_a\to\dots\to t_b$ be such a run. Its vertices $t_{a+1}>\dots>t_b$ are not in~$\Gamma$, so each lies in an
interval $[nP,nP+P-\mu)$ with $n\ge0$. If $t_i$ lies in such an interval and $t_{i+1}<nP$ for consecutive
vertices of the run, then $n\ge1$ and $t_{i+1}\ge t_i-\mu\ge nP-\mu$, so $t_{i+1}\bmod P\ge P-\mu$ and
$t_{i+1}\in\Gamma$, which is impossible. Hence $t_{a+1},\dots,t_b$ are distinct elements of one interval of $P-\mu$
consecutive integers, and the run has $b-a\le P-\mu$ edges. So $m\le D'+(D'+1)(P-\mu)$, and this is at most
$(D'+1)P-1$ as $\mu\ge1$. Each of the $\lceil T/P\rceil$ periods $[nP,(n+1)P)$ that meet $[0,T)$ contains $\mu$
revealed residues; for $\mu=1$ the revealed steps are $P-1,2P-1,\dots$, of which $\lfloor T/P\rfloor$ are less
than~$T$.
\end{informalproof}

\begin{remark}[Why blocks of $\mu$ steps]\label{rem:dec-jitter}
For $\mu=1$ one revealed step per period suffices, because a run of edges of length~$1$ visits consecutive times.
Single revealed steps do not cut runs of edges of length at most~$\mu$ when $\mu\ge2$: a work head that oscillates
between two adjacent cells has $\prev_j(t)=t-2$ for all $t\ge2$, and if $P$ is even, the path $t\to t-2\to t-4\to
\cdots$ through even times avoids every step $t\equiv P-1\pmod P$. A depth-reduction statement that treats such short
revisits (jitter, or revisits within a window of $\mu$ steps) as free must therefore be combined with
\cref{lem:dec-longedge} for that~$\mu$, which reveals blocks of $\mu$ consecutive steps at $\mu$ times the cost.
Theorem~B needs only $\mu=1$, because it counts every hop of length at least~$2$ as long.
\end{remark}

\begin{corollary}[Corollary B: home-based oblivious machines reduce to the crux]\label{cor:B}
\notformalized
Let $k\ge1$ and $\beta\ge2$ be integers, $C\ge1$, $A>0$ and $0<\eps\le1$, and put
\[
  D=D(\eps)=2^{\lceil456k^3(1+A)\beta C/\eps\rceil}\cdot\lceil4\beta C/\eps\rceil=2^{O(k^3(1+A)\beta C/\eps)}.
\]
Suppose that $\Crux(\mathcal S_{k,C,\beta,D},\rho)$ holds for some $\rho>\eps$. Then for every machine~$M$ with $k$
work tapes and value width at most~$\beta$ there is~$d_0$ such that for no $d\ge d_0$ do the following three
conditions hold together at $L=d^2b(d)$: $M$ is oblivious at length~$L$ with halting time $T\le C(L+1)$; each
work-head walk~$x_j$ of~$M$ on $[0,T)$ has leg area at most $AT$ measured from some cell~$c_j$ (that is, the walk
$x_j-c_j$ satisfies~(A)); and $M$ transposes all $d\times d$ matrices with $b(d)$-bit entries. In particular, if an
oblivious machine with $k$ work tapes and value width at most~$\beta$, whose work heads are home-based at every input
length, transposes within $C(L+1)$ steps for infinitely many~$d$, then for every $\eps\in(0,1]$ the crux
$\Crux(\mathcal S_{k,C,\beta,D(\eps)},\rho)$, with $A=1$ in $D(\eps)$, fails for every rate $\rho>\eps$.
\end{corollary}

\begin{informalproof}
Put $D_1=2^{\lceil456k^3(1+A)\beta C/\eps\rceil}$ and $P=\lceil4\beta C/\eps\rceil$, so that $D=D_1P$; note that
$D_1\ge4^k$, because $456k^3(1+A)\beta C/\eps\ge2k$. Fix~$M$, let $d_0$ be given by
$\Crux(\mathcal S_{k,C,\beta,D},\rho)$ for $\eta=(\rho-\eps)/2$, and suppose that the three conditions hold at some
$d\ge d_0$. As in the proof of \cref{thm:D}, $M$ appends exactly $L$ output bits, and on every input~$w$ its output
bit~$\ell$ is $w_{\operatorname{tpos}(\ell)}$.

Apply \cref{thm:B}(a) with depth parameter~$D_1$ to the translated work-head walks $x_1-c_1,\dots,x_k-c_k$ on
$[0,T)$, whose walk DAG is that of $x_1,\dots,x_k$: some cut~$F$ of depth at most $D_1-1$ has
$\card F\le114k^3(1+A)T/\log_2D_1\le\eps T/(4\beta C)$. Let $\Gamma_0$ be the set of
steps~$s$ such that some edge of~$F$ ends at~$s$; so $\card{\Gamma_0}\le\card F$. Every edge of a $\Gamma_0$-free path
in $H_M$ ends outside~$\Gamma_0$, so its endpoints are not those of an edge of~$F$. Read in the walk DAG, the path
therefore avoids~$F$, and its edges of length at least~$2$ are long hops; so it has at most $D_1-1$ edges of length
greater than~$1$. By \cref{lem:dec-longedge} with $\mu=1$, the genie
$\Gamma=\Gamma_0\cup\{t<T:t\equiv P-1\pmod P\}$ has depth at most $D_1P-1<D$, and
\[
  \beta\card\Gamma\;\le\;\beta\card F+\beta\lfloor T/P\rfloor\;\le\;\frac{\eps T}{4C}+\frac{\eps T}{4C}
  \;\le\;\frac{\eps(L+1)}2\;\le\;\eps L,
\]
using $T\le C(L+1)$ and $L\ge1$. So $I_\Gamma\in\mathcal S_{k,C,\beta,D}$, and by \cref{cor:dec-output} the revealed
word $G_\Gamma$ is an index code for $I_\Gamma$ with at most $\eps L$ bits, while the crux requires at least
$(\rho-\eta)L>\eps L$ bits. This contradiction proves the first statement. For the second, a work head that is
home-based with home cell~$c_j$ satisfies~(A) with $A=1$ when measured from~$c_j$, by \cref{lem:dec-area} applied to
$x_j-c_j$; a machine that transposes within $C(L+1)$ steps at length~$L$ has halting time at most $C(L+1)$ there; so a
machine as described would violate the first statement at infinitely many~$d$ if the crux held at some rate
$\rho>\eps$.
\end{informalproof}

The genie~$\Gamma$ built in this proof does not use that $M$ transposes. It shows that the conclusion of
$\mathrm{depth}\text{-}\TCP(k,C,\beta,\eps,D(\eps))$ holds for every machine~$M$ with $k$ work tapes and value width at
most~$\beta$, at every length $L=d^2b(d)$ at which $M$ is oblivious with halting time at most $C(L+1)$ and its
work-head walks satisfy~(A), measured from some home cells.

\paragraph{Why it matters.} Theorem~D needs two hypotheses. For home-based heads Theorem~B supplies the first, so
Corollary~B reduces home-based oblivious machines to the crux alone, whatever the coupling of their heads. Read as a
requirement on the crux, it suffices to prove $\Crux(\mathcal S_{k,C,\beta,D},\rho)$ at a single depth $D=D(\eps)$
with a rate $\rho>\eps$, that is, a rate exceeding about $456k^3(1+A)\beta C/\log_2D$ at depth~$D$ (\cref{rem:dec-rates});
this is far beyond the rate, at most $(k+1)^{-D-1}$, that the counting bound of \cref{prop:dec-counting} gives, but, for large~$D$, far below
rate~$\tfrac12$. Read as a design constraint, an oblivious linear-time transposer must either have a work head whose
leg area, measured from any home cell, is not $O(T)$, or produce, for every $\eps\in(0,1]$ and at every size at which it transposes, an index
code with at most $\eps L$ bits for its own cone side information at depth $D(\eps)$, and hence a storage code of
rate at least $1-\eps$ (\cref{rem:dec-crux}(b)).

\begin{remark}[Adaptive machines]\label{rem:dec-adaptive}
The argument uses obliviousness twice: the dependency DAG and the times of the output bits must be the same for all
inputs. We do not treat adaptive machines here. We expect that an extension of Corollary~B to adaptive machines whose heads are
home-based on every input would need additional hypotheses on head motion far from home and a stronger form of the
crux, for index codes that need to be correct only on a large set of inputs. No result of this paper depends on such an extension.
\end{remark}
